% Source: Lecture7.pdf, pages 1-3 (1D Ising model, continued). Page 8 is blank.
% The Potts-model half of the lecture (pages 3-7) is in lec07-potts.tex.
\lecture{7}
\subsubsection*{Correlation length and susceptibility of the 1D Ising model}

Previously, we calculated (for $H = 0$) the 2-spin correlation function
\begin{equation}
  \ev{\sigma_0\sigma_\ell} = v^\ell,\qquad \ell\ge 0,
\end{equation}
where $v = \tanh K$, $K = \beta J$. Recall that for $T>0$, since $|\tanh K|<1$,
\begin{enumerate}[label=(\arabic*)]
  \item $|\ev{\sigma_0\sigma_\ell}|$ is a monotonically decreasing function of $\ell\ge 0$, where $\ell$ is the distance
    between the spins in lattice units (if $\ell = 0$, $\ev{\sigma_0^2} = 1$);
  \item for fixed $\ell\ge 1$, $|\ev{\sigma_0\sigma_\ell}|$ is a monotonically decreasing function of $T$.
\end{enumerate}
These properties are physically understandable. Property (1) expresses the fact that as the two spins get farther and
farther apart, the value of spin $\sigma_0$ has a progressively smaller and smaller influence on the value of spin
$\sigma_\ell$. Property (2) follows because the higher the temperature $T$ is, i.e.\ the smaller $\beta$ and $|K|$ are, the
less is the ordering effect of the interaction term in $e^{-\beta\Ham} = e^{K\sum_{\langle ij\rangle}\sigma_i\sigma_j}$.

Also, note that the property that $\tanh(K)$ is an odd function of $K$ means that this correlation function has the
properties expected from the ferromagnetic (FM) and antiferromagnetic (AFM) signs of the spin--spin coupling, respectively
$J>0$ and $J<0$, namely that $\ev{\sigma_0\sigma_\ell}>0$ for FM $J$, while for AFM $J<0$,
\begin{equation}
  \ev{\sigma_0\sigma_\ell} = (\tanh K)^\ell = (-1)^\ell|\tanh K|^\ell
  = \begin{cases} >0 & \text{if $\ell$ even}\\ <0 & \text{if $\ell$ odd.}\end{cases}
\end{equation}
\begin{equation*}
  \text{FM order:}\quad
  \begin{array}{cccc} + & + & + & +\\[-2pt] \scriptstyle 0 & \scriptstyle 1 & \scriptstyle 2 & \scriptstyle 3\end{array}
  \ \text{or}\
  \begin{array}{cccc} - & - & - & -\\[-2pt] \scriptstyle 0 & \scriptstyle 1 & \scriptstyle 2 & \scriptstyle 3\end{array}
  \qquad\qquad
  \text{AFM order:}\quad
  \begin{array}{cccc} + & - & + & -\\[-2pt] \scriptstyle 0 & \scriptstyle 1 & \scriptstyle 2 & \scriptstyle 3\end{array}
  \ \text{or}\
  \begin{array}{cccc} - & + & - & +\\[-2pt] \scriptstyle 0 & \scriptstyle 1 & \scriptstyle 2 & \scriptstyle 3\end{array}
\end{equation*}

\paragraph{Correlation length.}

Next, calculate the correlation length $\xi$. Because of the bipartite nature of the 1D lattice, we may take $J>0$ with no
loss of generality. Recall that for $T\ne T_c$, a connected correlation function decays exponentially as the distance
between the spins gets large; $\ev{\sigma_0\sigma_r}\sim e^{-r/\xi}$ (in higher dimensions there may be power-law
prefactors, but here these are absent). So we have (with $v = \tanh K$)
\begin{equation}
  v^r = e^{-r/\xi}\ \Rightarrow\ v = e^{-1/\xi}\ \Rightarrow\ \ln v = -\frac{1}{\xi},\quad\text{i.e.}\quad
  \xi = -\frac{1}{\ln v} = \frac{1}{\ln\left(1/v\right)} .
\end{equation}
Since $v$ decreases with increasing $T$, $\ln(1/v)$ increases, and $\xi$ decreases. Again, this makes sense physically; as
$T$ increases, there is greater thermal disorder, so the correlation length should decrease.

Next, calculate the singularity in $\xi$ as $T\to 0$, i.e.\ $K\to\infty$:
\begin{align}
  \xi &= \frac{1}{\ln\left(1/v\right)};\qquad
  \frac{1}{v} = \frac{1}{\tanh K} = \frac{\cosh K}{\sinh K} = \frac{e^K+e^{-K}}{e^K-e^{-K}} = \frac{1+e^{-2K}}{1-e^{-2K}}\nonumber\\
  &\simeq\left(1+e^{-2K}\right)\left(1+e^{-2K}+\cdots\right) = 1+2e^{-2K},
\end{align}
so
\begin{equation}
  \ln\left(\frac{1}{v}\right)\simeq\ln\left(1+2e^{-2K}\right)\simeq 2e^{-2K}\ \Rightarrow\
  \xi\simeq\frac{1}{2e^{-2K}} = \frac{1}{2}e^{2K}\quad\text{as } T\to 0 .
\end{equation}
So the correlation length $\xi$ diverges exponentially rapidly as $T\to 0$,
\begin{equation}
  \xi\simeq\frac{1}{2}e^{2J/(\kB T)},
\end{equation}
an essential singularity. This is similar to what we found for $\bar\chi$, the initial ($H = 0$) magnetic susceptibility,
$\bar\chi_{H=0} = \beta e^{2K}$.

\paragraph{Susceptibility from the correlation function.}

It is useful to exhibit a different way to calculate $\bar\chi_{H=0}$. To do this we recall that, just as
$\bar M = \partial f/\partial h$ ``pulls down'' one $\sigma_i$ from $e^{-\beta\Ham}$,
\begin{equation}
  \bar M = \lim_{N\to\infty}\frac{1}{N}\,\frac{\sum_{\{\sigma_i\}}\left(\sum_j\sigma_j\right)e^{-\beta\Ham}}{\sum_{\{\sigma_i\}}e^{-\beta\Ham}},
\end{equation}
the susceptibility
\begin{equation}
  \bar\chi = \pder{\bar M}{H} = \beta\pder{\bar M}{h} = \beta\frac{\partial^2f}{\partial h^2}
\end{equation}
``pulls down'' two $\sigma$'s. For $T>0$, $\ev\sigma = \bar M = 0$, and
\begin{equation}
  \beta^{-1}\bar\chi = \lim_{N\to\infty}\frac{1}{N}\sum_{\text{all }i,j}\ev{\sigma_i\sigma_j},
\end{equation}
where, as indicated, $i$ and $j$ range over all sites. Since the lattice is homogeneous, this
$\lim_{N\to\infty}\sum_{i,j}$ double sum is equal to the sum over one site with the other fixed:
\begin{equation}
  \lim_{N\to\infty}\frac{1}{N}\sum_{\text{all }i,j}\ev{\sigma_i\sigma_j} = \sum_j\ev{\sigma_0\sigma_j},\qquad j\in\mathbb Z .
\end{equation}
Further, $\ev{\sigma_0\sigma_j} = \ev{\sigma_0\sigma_{-j}} = v^{|j|}$. So
\begin{equation}
  \sum_{j\in\mathbb Z}\ev{\sigma_0\sigma_j} = \underbrace{\ev{\sigma_0\sigma_0}}_{1}
  + \sum_{j=-\infty}^{-1}\ev{\sigma_0\sigma_j} + \sum_{j=1}^\infty\ev{\sigma_0\sigma_j}
  = 1 + 2\sum_{j=1}^\infty\ev{\sigma_0\sigma_j} = 1 + 2\sum_{j=1}^\infty v^j .
\end{equation}
Now $\sum_{j=0}^\infty v^j = 1/(1-v)$, provided this series converges, which it does for $|v|<1$, so
\begin{equation}
  \sum_{j=1}^\infty v^j = -v^0 + v^0 + \sum_{j=1}^\infty v^j = -1 + \sum_{j=0}^\infty v^j = -1 + \frac{1}{1-v} = \frac{v}{1-v}.
\end{equation}
So $1 + 2\sum_{j=1}^\infty v^j = 1 + 2\left(v/(1-v)\right) = (1+v)/(1-v)$, so
\begin{align}
  \beta^{-1}\bar\chi &= \frac{1+v}{1-v} = \frac{1+\tanh K}{1-\tanh K} = \frac{\cosh K+\sinh K}{\cosh K-\sinh K}\nonumber\\
  &= \frac{\left(e^K+e^{-K}\right)+\left(e^K-e^{-K}\right)}{\left(e^K+e^{-K}\right)-\left(e^K-e^{-K}\right)} = \frac{2e^K}{2e^{-K}} = e^{2K},
\end{align}
i.e.\ $\bar\chi = \beta e^{2K}$, as we derived earlier.

\paragraph{Susceptibility in a nonzero field.}

Next, discuss our earlier result in the case of a nonzero external magnetic field $H$. Recall the expression
\begin{equation}
  \bar\chi = \frac{\beta e^{-4K}\cosh h}{\left[\sinh^2h+e^{-4K}\right]^{3/2}}.
\end{equation}
Consider the behavior as $T\to\infty$, i.e.\ $\beta\to 0$, $K\to 0$ and $h\equiv\beta H\to 0$, so $\sinh h\to 0$,
$e^{-4K}\to 1$, $\cosh h\to 1$, $\bar\chi\to\beta$. This is the realization, in this model, of the Curie (Pierre Curie)
``law'', that $\bar\chi\propto\beta\propto1/T$ at high temperature $T$. This is observed experimentally in 3D as well as
lower-dimensional substances.

Consider the behavior as $T\to 0$. Here, with $H\ne 0$, one must analyze the $J>0$ (FM) and $J<0$ (AFM) cases separately.
This is clear, since an external magnetic field $H$ biases the system toward FM spin ordering ($h = \beta H$).
\begin{enumerate}[label=(\arabic*)]
  \item Case $J>0$: here, as $T\to 0$, $K\to\infty$ so $e^{-4K}\to 0$, and $h\to\infty$:
    $\cosh h\to e^h/2$, $\sinh h\to e^h/2$, $\cosh h/(\sinh^3h)\to(e^h/2)/(e^h/2)^3 = 4e^{-2h}$, so
    \begin{equation}
      \bar\chi\to\frac{\beta e^{-4K}\cosh h}{\sinh^3h}\to 4\beta e^{-4K}e^{-2h}\to 0 .
    \end{equation}
    This vanishing of $\bar\chi$ for $H\ne 0$ is physically understandable, because the external magnetic field biases the
    spins toward FM order, and this approaches a complete ordering as $T\to 0$, $|\bar M|\to 1$, so increasing $H$ has a
    negligible effect and $\bar\chi = \partial \bar M/\partial H\to 0$.
  \item Case $J<0$: here one has competing interactions: (i) the spin--spin interaction prefers the spins to order
    antiferromagnetically: $+-+-+-$; (ii) the external magnetic field prefers the spins to order ferromagnetically. This is
    an exactly solvable example of competing interactions. Which one dominates over the other? This will be a homework
    problem.
\end{enumerate}
